% Block 2 of the slides (02_agents.tex).
\section{Chatbots, harnesses and agents}\label{sec:agents}

\onslides{a chatbot where you are the loop; an agent where the loop closes; one
agent transcript; chatbot versus agent table; permissions; the live demo ``Same task, two
ways''; asking for verifiable output beyond code.}

\subsection{What a harness is}

The model of Section~\ref{sec:llm} maps text to text and does nothing else. Everything
else in a product such as Claude Code or Codex is the \emph{harness}: the program that
\begin{enumerate}
  \item assembles the context (system instructions, instruction files, your messages,
        earlier tool results),
  \item offers the model a set of \emph{tools}, each described by a name, a description
        and a schema for its arguments (read a file, edit a file, run a shell command,
        search the web),
  \item recognises when the model's output is a \emph{tool call}, a structured request
        such as \texttt{run("julia --project -e 'using Pkg; Pkg.test()'")}, asks you for
        permission if needed, executes it, and appends the result to the context,
  \item loops until the model answers without a tool call, and
  \item manages the context when it fills up (Section~\ref{sec:context}).
\end{enumerate}
The model never touches your computer. It writes text that requests an action, and the
harness decides whether to carry it out. A chatbot in the browser is a harness with no
tools on your machine, and \emph{you} carry code and error messages between the model and
your computer: you are the loop.

Interleaving reasoning with actions and observations is known as the ReAct pattern
\rf{react}. Schluntz and Zhang \rf{agents} distinguish \emph{workflows}, in which code
fixes the sequence of model calls, from \emph{agents}, in which the model decides its next
step from the results of the previous one. Their advice, start with the simplest design
that works, applies to how you use agents too.

\subsection{Why the closed loop matters}

The transcript on the slide is illustrative (the real test is the $\pi$-pulse line of the
\cmd{"Rabi"} testset, Section~\ref{sec:structure}), and it shows the point of the block. The agent runs the tests, reads a failure (\lstinline|0.0 ≈ 1.0|), reads the
source, finds $\Omega$ where $\Omega/2$ belongs, edits, and runs the tests again. It
found the factor of two \emph{because a test encoded the physics}: a resonant
$\pi$-pulse must invert the spin. With the bug, the population at $t = \pi/\Omega$ is
$\sin^2(\pi) = 0$ instead of $\sin^2(\pi/2) = 1$ (Section~\ref{sec:failures}). An agent
can check its work against reality, but only against the reality you give it.

Evidence that this works at scale: SWE-bench \rf{swebench} tests whether a system can
resolve 2294 real GitHub issues from 12 popular Python repositories, judged by the repositories'
own tests. METR measures the length of software tasks, in human working time, that
agents complete with $50\,\%$ reliability, and found that this length has doubled
roughly every seven months since 2019 \rf{metr-horizon}. Both benchmarks rely on
tests written by people.

\subsection{Permissions}

Agents ask before editing a file or running a command, and every agent lets you approve
whole classes of actions (all edits, or specific commands). The slide's rules:
\begin{itemize}
  \item Never auto-approve deletion, overwriting, or anything outside the project
        folder. An agent runs commands with your user's rights.
  \item ``Yes, yes, yes'' is not supervision. If you approve without reading, you have
        granted full permission with extra steps.
  \item Git makes a longer leash safe (Section~\ref{sec:git}): any change inside the
        repository can be inspected and undone, but nothing outside it can.
\end{itemize}

\subsection{The live demo: same task, two ways}

The prompt on the slide asks for Rabi oscillations with the course Hamiltonian, two
detunings, three Rabi periods and a check against the analytic formula. What to point out
while it runs:
\begin{itemize}
  \item In the chat, every error costs a copy--paste round trip; the model never sees the
        plot.
  \item The agent writes files, runs Julia, reads the error, fixes it. Count the tool
        calls.
  \item Ask the room: what did it check? (Typically it compares with the formula at a
        few points.) What did it not check? (Units, the sign convention of $\Delta$, the
        initial state, that the plot shows what the caption claims.)
\end{itemize}

\subsection{Beyond code}

For any task, ask for an output you can verify: links you open for a literature search, a
script that checks a derivation numerically, an analysis that lives as a script in the
repository, and text you can defend line by line. A number in a chat window cannot be
reproduced; a script in a repository can.

\begin{qa}
\Q{Which agent should I use?}
The course does not prescribe one; the concepts carry over. Use the one your institution
or subscription gives you, and check its data-handling terms (Section~\ref{sec:failures}).

\Q{Can the agent delete my files?}
If you permit it, yes. Keep the project under git, commit before the agent starts, and do
not approve commands that act outside the project folder.

\Q{Is ``agent'' just marketing?}
In the sense used here it has a precise meaning: a model in a loop that chooses tool
calls from their results \rf{agents}. Whether that is useful depends on whether the
results contain a signal of correctness, which is what tests provide.
\end{qa}

\textbf{Read next:} \emph{Building effective agents} \rf{agents} (short, practical);
Anthropic's best-practice guide for its coding agent \rf{cc-best}, whose first advice,
``give Claude a way to verify its work'', is tool-specific in wording and general in
content.
